%
% definition.tex -- Basics
%
% (c) 2020 Prof Dr Andreas Müller, Hochschule Rapperswil
%
% !TEX root = ../../buch.tex
% !TEX encoding = UTF-8
%

\section{Algebraische Definition\label{qa:section:definition}}
\kopfrechts{Algebraische Definition}

Naheliegend wäre, für dreidimensionale Drehungen auch dreidimensionale
Zahlen zu verwenden. Mit $\mathbb{R}, i, j$ hätte man genau so viele, wie der
Anschauungsraum besitzt. An diesem Ansatz ist Hamilton gescheitert, 
und zwar an der Multiplikation. Es werden vier Dimensionen benötigt, 
welche Schritt für Schritt erläutert werden.


\subsection{Addition und Multiplikation komplexer Zahlen}
Sei eine positive reelle Zahl $x$: Wenn $x$ mit $i$ multipliziert 
wird, dann entsteht aus einer rein reeller Zahl eine rein imaginäre 
Zahl $ix$. Dies kann als eine 90° Drehung im Gegenuhrzeigersinn 
interpretiert werden. Wird nun das Ergebnis $ix$ ein weiteres Mal 
mit $i$ multipliziert, sodass $iix = i^2x$, dann greift die oben 
genannte Regel und $-x$ entsteht. Die Zahl wurde somit um weitere 
90° im Gegenuhrzeigersinn gedreht. Die selbe Regel gilt auch 
allgemein: Selbst wenn $x$ keine reelle positive Zahl, sondern eine 
komplexe Zahl ist, geschieht die Drehung auf gleiche Weise. 

Für eine ganze 360° Drehung muss $x$ vier mal mit $i$ multipliziert 
werden. Potenzen mit Basis $i$ wiederholen sich also nach vier 
Iterationen.

% Tabelle und Bild einfügen

Mit Multiplikationen können nicht nur Drehungen vollzogen werden., 
Statt eine Zahl mit $i$ zu multiplizieren, kann mit einer 
Multiplikation von zwei komplexen Zahlen $c_1, c_2$ eine präzise 
Platzierung mit bestimmter Drehung und Skalierung durchgeführt 
werden. Dies geschieht distributiv, wobei das Ergebnis mit
\begin{equation*}
    (a + ib) \cdot (c + id) 
    = ac + iad + ibc - bd = 
    (ac - bd) + i(ad + bc), \qquad a,b,c,d \in \mathbb{R}
\end{equation*}
wieder in Real- und Imaginärteil aufgeteilt wird.

Die Addition zweier komplexer Zahlen gleicht dem Vorgang einer 
Vektoraddition, bei der die Zahlen komponentenweise zusammengezählt
werden, was zu
\begin{equation*}
    (a + ib) + (d + id) = a+c + ib + id = (a+c) + i(b+d)
\end{equation*}
führt.
Auch hier wird das Ergebnis auf den Real- und Imaginärteil 
aufgeteilt. Quaternionen verhalten sich sehr ähnlich, wie im 
nächsten Abschnitt beschrieben ist.

\subsection{Addition und Multiplikationen von Quaternionen}
Für Quaternionen werden zwei neue Regeln eingeführt. Die erste 
Regel ist die selbe wie bei den komplexen Zahlen. Sie besagt, dass 
das Quadrat einer imaginären Einheit jeweils $-1$ ergibt, was sich in
\begin{equation*}
    i^2 = j^2 = k^2 = -1
\end{equation*}
zusammenfassen lässt.
Die zweite Regel führt eine Abhängigkeit zwischen den imaginären 
Einheiten ein. Eine Folge der drei Imaginärachsen in alphabetischer 
Reihenfolge $i\!jk$ ergibt ebenfalls $-1$. Dies gilt auch für 
zyklische Verschiebungen, wie in
\begin{equation*}
    i\!jk = ki\!j = jki = -1
\end{equation*}
beschrieben. Damit entstehen weitere Eigenschaften: Da $i\!jk = kk = 
-1$, muss $i\!j = k$ entsprechen. Dasselbe gilt auch für $ki = j$ und 
$jk = i$. Die Reihenfolge einer Multiplikation spielt dabei eine 
grosse Rolle, denn $ik \neq ki$. Dies kann mit dem folgenden
Experiment
\begin{align*}
    ik &= i \cdot 1 \cdot k \\
       &= -i \cdot j^2 \cdot k \\
       &= -\underbrace{i \cdot j}_k \cdot 
    \underbrace{j \cdot k}_{i} \\
        \Rightarrow ik &= -ki
\end{align*}
nachgewiesen werden.
Wenn die Reihenfolge der Multiplikation vertauscht wird, dann 
ändert sich das Vorzeichen. Quaternionen sind also nicht 
kommutativ. Dies kann intuitiv verstanden werden, wenn man sich 
\index{Quaternionen!nicht kommutativ}%
eine Drehung um die $x$-Achse und eine Drehung um die $y$-Achse 
vorstellt. Die Reihenfolge der beiden Drehungen spielt eine Rolle, 
da das Ergebnis unterschiedlich ausfällt.

Die Addition von zwei Quaternionen funktioniert analog zu den 
komplexen Zahlen.
\index{Quaternionen!Addition}%
Die Komponenten werden auch hier komponentenweise
addiert, sodass
\begin{equation*}
    (a + ib + jc + kd) + (e + if + jg + kh) 
    = (a+e) + i(b+f) + j(c+g) + k(d+h)
\end{equation*}
gilt.
Auch das Multiplizieren zweier Quaternionen funktioniert analog zu 
den komplexen Zahlen. Die Komponenten werden distributiv
\index{Quaternionen!Multiplikation}%
multipliziert und das Ergebnis wieder in Real- und Imaginärteile
aufgeteilt, woraus
%\begin{align*}
%    (a + ib + jc + kd) \cdot (e + if + jg + kh) 
%    &= \qquad ae + aif + ajg + akh \\
%    &\qquad+ ibe + i^2bf + i\!jbg + ikbh \\
%    &\qquad+ jce + jif + j^2cg + jkch \\
%    &\qquad+ kde + kif + kjcg + k^2dh \\
%    &= \qquad (ae - bf - cg - dh) \\
%    &\qquad+ i(af + be + ch - dg) \\
%    &\qquad+ j(ag - bh + ce + df) \\
%    &\qquad+ k(ah + bg - cf + de)
%\end{align*}
\begin{equation*}
\renewcommand{\arraycolsep}{1.5pt}
\renewcommand{\arraystretch}{1.2}
\begin{array}{rccrcrcrcr}
(a + ib + jc + kd) \cdot (e + if + jg + kh) 
    &=& &  ae                  &+& a{\color{darkred}i}f   &+& a{\color{darkred}j}g    &+& a{\color{darkred}k}h \\
    & &+& {\color{darkred}i}be &+& {\color{darkred}i^2}bf &+& {\color{darkred}i\!j}bg &+& {\color{darkred}i}kbh \\
    & &+& {\color{darkred}j}ce &+& {\color{darkred}ji}cf  &+& {\color{darkred}j^2}cg  &+& {\color{darkred}jk}ch \\
    & &+& {\color{darkred}k}de &+& {\color{darkred}ki}df  &+& {\color{darkred}k\!j}dg   &+& {\color{darkred}k^2}dh \\[5pt]
    &=& &\rlap{\hspace*{-9pt}$\,(ae {\color{darkred}\mathstrut-\mathstrut} bf {\color{darkred}\mathstrut-\mathstrut} cg {\color{darkred}\mathstrut-\mathstrut} dh)$}   \\
    & &+&\rlap{\hspace*{-9pt}$ \llap{${\color{darkred}i}$}\,(af + be {\color{darkred}\mathstrut+\mathstrut} ch {\color{darkred}\mathstrut-\mathstrut} dg)$} \\
    & &+&\rlap{\hspace*{-9pt}$ \llap{${\color{darkred}j}$}\,(ag {\color{darkred}\mathstrut-\mathstrut} bh + ce {\color{darkred}\mathstrut+\mathstrut} df)$} \\
    & &+&\rlap{\hspace*{-9pt}$ \llap{${\color{darkred}k}$}\,(ah {\color{darkred}\mathstrut+\mathstrut} bg {\color{darkred}\mathstrut-\mathstrut} cf + de)$}
\end{array}
\end{equation*}
folgt.
Farbige Operationszeichen rühren von Produkten zweier Basisquaternionen her.

\subsection{Warum vier Dimensionen?}
Angenommen, es gäbe eine dreidimensionale Zahlenmenge mit Basis $1, i, j$ und
$i^2 = j^2 = -1$, in der uneingeschränkt dividiert werden kann. Dann
müsste insbesondere das Produkt $i\!j$ wieder in dieser Menge liegen,
sich also als Linearkombination
\begin{equation*}
    i\!j = a + bi + cj, \qquad a,b,c \in \mathbb{R}
\end{equation*}
schreiben lassen. Multipliziert man diese Gleichung von links mit $i$,
so ergibt die Gleichung $i(i\!j) = i^2 j = -j$, die Linearkombination
dagegen
\begin{align*}
    i(a + bi + cj) &= ai + bi^2 + c(i\!j) \\
                   &= -b + ai + c(a + bi + cj) \\
                   &= (ca - b) + (a + cb)i + c^2\!j .
\end{align*}
Ein Koeffizientenvergleich beim Term $j$ liefert $c^2 = -1$, was für
eine reelle Zahl $c$ unmöglich ist. Das Produkt $i\!j$ lässt sich also
nicht im dreidimensionalen Raum unterbringen.

Hamiltons Ausweg bestand darin, dieses Produkt nicht mehr erklären zu
wollen, sondern es als neue, vierte Basisgrösse zu akzeptieren:
$k := i\!j$. Die Erweiterung ist damit zwangsläufig vierdimensional. Der
Preis dafür ist die Kommutativität, denn aus $k = i\!j$ folgt $ji = -k$,
wie oben bereits gezeigt wurde.

\subsection{Konjugation, Norm und Inverse}

Auch für Quaternionen existiert die \emph{komplexe Konjugation}. Wie bei 
\index{Quaternionen!Konjugation}%
den komplexen Zahlen wird der Realteil nicht verändert, während die
Imaginärteile negiert werden, also
\begin{equation*}
    q^* = (a + ib + jc + kd)^* = a - ib - jc - kd .
\end{equation*}
Die euklidische \emph{Norm} einer Quaternion ist die Summe der Quadrate aller 
Komponenten. Sie wird durch
\begin{equation*}
    |q|^2 = |a + ib + jc + kd|^2 = a^2 + b^2 + c^2 + d^2
\end{equation*}
berechnet.
Die Norm lässt sich auch über die Konjugation als
\index{Quaternionen!Norm}%
\begin{equation*}
    |q|^2 = q \cdot q^* 
    = (a + ib + jc + kd) \cdot (a - ib - jc - kd)
\end{equation*}
schreiben.
\emph{Einheitsquaternionen} sind Quaternionen mit der Norm 1. Sie 
\index{Quaternionen!Einheitsquaternion}%
bilden die Einheitskugel im vierdimensionalen Raum, wie die komplexen 
Zahlen z mit $|z| = 1$ den 
\index{Quaternionen!Einheitsquaternionen}%
Einheitskreis im komplexen Raum. In den folgenden Abschnitten wird 
gezeigt, dass vor allem Einheitsquaternionen von grosser 
Wichtigkeit sind.

Um die \emph{Inverse} einer Quaternion zu berechnen, werden die Norm und 
die Konjugation benötigt. Das Ziel ist, dass $q q^{-1} = 1$ 
\index{Quaternionen!Inverse}%
ergibt. Diese Gleichung lässt sich zu
\begin{align*}
    q q^{-1} &= 1 \\
    q^* (qq^{-1}) &= q^* \\
    (q^* q) q^{-1} &= q^* \\
    \underbrace{|q|^2}_{\in \mathbb{R}} q^{-1} &= q^* \\
    q^{-1} &= \frac{q^*}{|q|^2}
\end{align*}
umstellen.
Was auffällt ist, dass die Inverse einer Einheitsquaternion gleich 
der Konjugation ist. Damit lassen sich Einheitsquaternionen sehr 
einfach invertieren.
